% P33の方法節部品。独自方式の提案・結果・投稿本文ではない。
% モデルと予算はconfigs/P33_output_views_v01.jsonを基準とする。
\section{出力省略の比較}
\subsection{目的と共通状態}
AIハーネスによる入力削減は，その後の再取得や追加生成を含む総費用と，
最終成果の品質によって評価する必要がある．
本比較では，実行コードと作業状態を保ったまま古い出力を省く通常方式を調べる．
出力省略や変数一覧の機能そのものを提案技術とはしない．

開発用の模擬アプリ操作４課題について，記録済みの実行の第８応答後を
共通の分岐点とする．課題は写真整理，共有・依頼，返金，曲一覧の書出である．
各課題についてアプリの状態，Python変数，課題要求，公開済みAPI，
原記録の参照手段，モデルと推論設定，後続予算を共通にする．
状態の再構成は公開返却・保存DB・プログラムの再生結果で照合し，
全プロセスの同一性を保証したとはしない．

\subsection{比較条件}
条件Fは全ての公開会話履歴を渡す．条件Mは最後の完全な１ターンを残し，
それより古い成功したPython標準出力を，原記録IDを含む省略表示へ置換する．
実行コード，結果状態，API受領記録，例外，切り詰めの有無は保持する．
未確定の操作結果や完了宣言を省略対象にせず，
置換後の表現が元より長くなる短い出力も保持する．
原記録は削除せず，全条件から同じ方法で再読できる．
省略は分岐時に一度だけ適用し，以後の会話は各条件に追加する．

条件MIは条件Mへ，境界時点のPython変数の名前・型・要素数と，
公開済みAPIの名前一覧を加える．値そのものや，未公開のAPIは追加しない．
一覧の作成費用と入力トークンを計上する．
全条件へ共通の省略表示と再読方法の説明を渡す．
このためMとMIの差には一覧の内容とその追加入力が含まれる．

全履歴の対照も新しく実行し，過去の最良試行を選ばない．
４課題と３条件の計12試行を，条件順を課題ごとに入れ替えて直列に実行する．
各条件は１試行であり，12独立課題とは数えない．
後続上限は16生成，総入力・出力120,000トークン，480秒とする．
前半のAPI回数を課題上限へ含め，実験用の環境再構成は別途記録する．

\subsection{判定の範囲}
最終成果を独立に評価し，課題ごとの成否と全費用を保存する．
開発上の採否は，対照で成功した課題の成功を失わず，
４課題合計の総トークンが減り，両条件に使用量不明がない場合を通過とする．
別の課題の成功で悪化を相殺しない．
この少数開発比較を品質の非劣性や未経験課題への一般化の証明とはしない．

生成回数を$n$，初回入力を$I$，各回入力を$I_j$，全出力を$O$とすると，
\begin{equation}
 T=nI+\sum_j(I_j-I)+O
\end{equation}
により，初回の短縮と後続の費用を分けて記述できる．
これは観測した使用量の恒等式であり，各項の独立した因果効果ではない．
キャッシュ利用，実時間，API回数は総トークンと別に示す．
